\documentclass[11pt]{article}

\usepackage{amsmath,amssymb,amsthm,mathtools}
\usepackage{hyperref}
\usepackage{enumitem}

\newtheorem{theorem}{Theorem}
\newtheorem{lemma}{Lemma}
\newtheorem{proposition}{Proposition}
\newtheorem{corollary}{Corollary}
\theoremstyle{definition}
\newtheorem{definition}{Definition}
\newtheorem{remark}{Remark}

\newcommand{\Bbbk}{\mathbf{k}}
\newcommand{\Uc}{\mathcal{U}}
\newcommand{\fsl}{\mathfrak{sl}}
\newcommand{\fb}{\mathfrak{b}}
\newcommand{\Id}{\mathrm{id}}
\newcommand{\eps}{\varepsilon}
\newcommand{\del}{\delta}
\newcommand{\bipr}{\bowtie}
\newcommand{\tensor}{\otimes}
\newcommand{\la}{\langle}
\newcommand{\ra}{\rangle}

\newcommand{\Khat}{\widehat{K}}
\newcommand{\E}{E}
\newcommand{\F}{F}

\DeclareMathOperator{\End}{End}

\begin{document}

\title{The punctured bigon algebra and the generalized Drinfeld double of the quantum Borel subalgebras of $\Uc_q(\fsl_n)$}
\author{}
\date{}
\maketitle

\begin{abstract}
Let $\Uc_q(\fb_\pm)$ be the (extended, i.e.\ simply connected) quantum Borel Hopf subalgebras of $\Uc_q(\fsl_n)$, and let $\la\cdot,\cdot\ra$ be the standard Drinfeld--Tanisaki skew Hopf pairing on $\Uc_q(\fb_+)\tensor \Uc_q(\fb_-)$. We construct an explicit algebra isomorphism between the ``punctured bigon algebra'' $\bcS_n(\mathfrak B^*)$ (presented by $RLL$--relations on two triangular matrices $(b_{ij})$ and $(c_{ij})$) and the generalized quantum double $\Uc_q(\fb_-)\bipr \Uc_q(\fb_+)$. We then show that $\Uc_q(\fsl_n)$ is the quotient of $\bcS_n(\mathfrak B^*)$ by the two-sided ideal generated by $b_{ii}c_{ii}-1$ for $1\le i\le n$.
\end{abstract}

\subsection*{Conventions and ground ring}
Fix an integer $n\ge 2$ and a field $\Bbbk$ of characteristic $0$ together with $q\in \Bbbk^\times$ such that $q$ is not a root of unity. (All constructions below make sense over $\mathbb{Q}(q)$ as well.) We use Sweedler notation $\Delta(x)=\sum x_{(1)}\tensor x_{(2)}$ and $\Delta^{(2)}(x)=\sum x_{(1)}\tensor x_{(2)}\tensor x_{(3)}$.

%%%%%%%%%%%%%%%%%%%%%%%%%%%%%%%%%%%%%%%%%%%%%%%%%%%%%%%%%%%%%%%%%%%%%%%%%%%%%%%

\section{Skew Hopf pairings and generalized quantum doubles}

\begin{definition}[Skew Hopf pairing]
Let $A,B$ be Hopf algebras over $\Bbbk$. A bilinear map $\la\cdot,\cdot\ra:A\tensor B\to \Bbbk$ is a \emph{skew Hopf pairing} if, for all $a,a'\in A$ and $b,b'\in B$,
\begin{enumerate}[label=(\roman*)]
\item $\la aa',b\ra=\sum \la a,b_{(2)}\ra\,\la a',b_{(1)}\ra$,
\item $\la a,bb'\ra=\sum \la a_{(1)},b'\ra\,\la a_{(2)},b\ra$,
\item $\la 1,b\ra=\eps(b)$ and $\la a,1\ra=\eps(a)$,
\item $\la S_A(a),b\ra=\la a,S_B^{-1}(b)\ra$.
\end{enumerate}
\end{definition}

\begin{definition}[Generalized quantum double as a double cross product]
Let $A,B$ be Hopf algebras equipped with a skew Hopf pairing $\la\cdot,\cdot\ra:A\tensor B\to \Bbbk$.
Define $\la\cdot,\cdot\ra^{-1}:A\tensor B\to \Bbbk$ by $\la a,b\ra^{-1}:=\la S_A(a),b\ra$ (equivalently $\la a,S_B^{-1}(b)\ra$).
The \emph{generalized quantum double} (or \emph{generalized Drinfeld double})
\[
A \bipr B
\]
is the Hopf algebra whose underlying coalgebra is $A\tensor B$ (tensor coalgebra), and whose multiplication is determined by
\begin{equation}\label{eq:double-mult}
(a\bipr b)\cdot (a'\bipr b')
=\sum 
\la a'_{(1)}, b_{(1)}\ra\,
\la a'_{(3)}, b_{(3)}\ra^{-1}\,
(aa'_{(2)}\bipr b_{(2)}b').
\end{equation}
Here we write $a\bipr b$ for $a\tensor b$ to emphasize the twisted product.
\end{definition}

\begin{remark}
Formula \eqref{eq:double-mult} is standard in the literature on generalized doubles; see e.g.\ Heckenberger--Schneider for the construction and its basic properties.  % ([arxiv.org](https://arxiv.org/abs/0911.2282))
\end{remark}

\begin{lemma}[Cross relation in the generalized double]\label{lem:cross-relation}
Inside $A\bipr B$, the canonical embeddings $A\hookrightarrow A\bipr B$, $a\mapsto a\bipr 1$ and $B\hookrightarrow A\bipr B$, $b\mapsto 1\bipr b$ are algebra maps. Moreover, for $a\in A$, $b\in B$ one has
\begin{equation}\label{eq:cross}
(1\bipr b)\,(a\bipr 1)
=
\sum 
\la a_{(1)}, b_{(1)}\ra\,
\la a_{(3)}, b_{(3)}\ra^{-1}\,
(a_{(2)}\bipr b_{(2)}).
\end{equation}
\end{lemma}

\begin{proof}
Set $a'=a$, $b=b$, $a=1$, $b'=1$ in \eqref{eq:double-mult} and use $\eps$ and $S$ axioms of a skew Hopf pairing.
\end{proof}

%%%%%%%%%%%%%%%%%%%%%%%%%%%%%%%%%%%%%%%%%%%%%%%%%%%%%%%%%%%%%%%%%%%%%%%%%%%%%%%

\section{The simply connected quantum group $\Uc_q(\fsl_n)$ and its Borel subalgebras}

We recall the extended (simply connected) Drinfeld--Jimbo quantum group, as in e.g.\ Klimyk--Schm\"udgen (where it is denoted $U_q^{ext}(\fsl_n)$).

\begin{definition}[Extended $\Uc_q(\fsl_n)$]
Let $\Uc_q(\fsl_n)$ be the Hopf algebra over $\Bbbk$ generated by
\[
\E_i,\F_i\quad (1\le i\le n-1), 
\qquad
\Khat_j^{\pm 1}\quad (1\le j\le n),
\]
subject to the relations:
\begin{enumerate}[label=(\roman*)]
\item $\Khat_j\Khat_j^{-1}=1$, $\Khat_j\Khat_k=\Khat_k\Khat_j$, and $\Khat_1\cdots \Khat_n=1$;
\item $\Khat_j \E_i \Khat_j^{-1}=q^{\del_{j,i}-\del_{j,i+1}}\E_i$ and $\Khat_j \F_i \Khat_j^{-1}=q^{-\del_{j,i}+\del_{j,i+1}}\F_i$;
\item $[\E_i,\F_j]=\del_{i,j}\,\dfrac{K_i-K_i^{-1}}{q-q^{-1}}$, where $K_i:=\Khat_i\Khat_{i+1}^{-1}$;
\item the $q$--Serre relations among the $\E_i$ and among the $\F_i$ (type $A_{n-1}$).
\end{enumerate}
The Hopf structure is given on generators by
\[
\Delta(\Khat_j)=\Khat_j\tensor \Khat_j,\quad
\Delta(\E_i)=\E_i\tensor 1 + K_i\tensor \E_i,\quad
\Delta(\F_i)=\F_i\tensor K_i^{-1}+1\tensor \F_i,
\]
together with the obvious counit and antipode.
\end{definition}

\begin{definition}[Quantum Borel Hopf subalgebras]
Let $\Uc_q(\fb_+)\subset \Uc_q(\fsl_n)$ (resp.\ $\Uc_q(\fb_-)\subset \Uc_q(\fsl_n)$) be the Hopf subalgebra generated by the $\E_i$ and the $\Khat_j^{\pm1}$ (resp.\ by the $\F_i$ and the $\Khat_j^{\pm1}$).
\end{definition}

%%%%%%%%%%%%%%%%%%%%%%%%%%%%%%%%%%%%%%%%%%%%%%%%%%%%%%%%%%%%%%%%%%%%%%%%%%%%%%%

\section{The Drinfeld--Tanisaki skew Hopf pairing}

\begin{theorem}[Drinfeld--Tanisaki pairing]\label{thm:DT-pairing}
There exists a unique skew Hopf pairing
\[
\la\cdot,\cdot\ra:\Uc_q(\fb_+)\tensor \Uc_q(\fb_-)\to \Bbbk
\]
such that, for $1\le i\le n-1$ and $1\le j,k\le n$,
\begin{equation}\label{eq:pairing-on-gens}
\la \Khat_j,\Khat_k\ra = q^{\frac1n-\del_{j,k}},\qquad
\la \E_i,\F_i\ra=\frac{1}{q^{-1}-q},\qquad
\la \E_i,\F_j\ra=0\ (i\ne j),
\end{equation}
and
\[
\la \Khat_j,\F_i\ra=\la \E_i,\Khat_j\ra=0.
\]
\end{theorem}

\begin{proof}
This is the standard Drinfeld pairing (in type $A_{n-1}$, extended/simply-connected normalization) due to Drinfeld and Tanisaki. A detailed construction and proof of the Hopf pairing axioms can be found in the standard references on quantum groups; see e.g.\ Klimyk--Schm\"udgen or Tanisaki's work on invariance of the Drinfeld pairing.
\end{proof}

%%%%%%%%%%%%%%%%%%%%%%%%%%%%%%%%%%%%%%%%%%%%%%%%%%%%%%%%%%%%%%%%%%%%%%%%%%%%%%%

\section{The generalized quantum double of the Borels}

\begin{definition}
Let
\[
\mathcal{D}:=\Uc_q(\fb_-)\bipr \Uc_q(\fb_+)
\]
be the generalized quantum double with respect to the Drinfeld--Tanisaki skew pairing of Theorem~\ref{thm:DT-pairing}.
We denote by
\[
\iota_-:\Uc_q(\fb_-)\hookrightarrow \mathcal{D},\qquad x\mapsto x\bipr 1,
\qquad
\iota_+:\Uc_q(\fb_+)\hookrightarrow \mathcal{D},\qquad y\mapsto 1\bipr y
\]
the canonical embeddings.
\end{definition}

\begin{remark}
It is a standard and very general fact that the (extended) quantized enveloping algebra can be recovered as a quotient of such a generalized double, by identifying the two copies of the Cartan torus; see e.g.\ Heckenberger--Schneider for the general ``$U_q$ as a quotient of a generalized quantum double'' statement.  % ([arxiv.org](https://arxiv.org/abs/0911.2282))
\end{remark}

%%%%%%%%%%%%%%%%%%%%%%%%%%%%%%%%%%%%%%%%%%%%%%%%%%%%%%%%%%%%%%%%%%%%%%%%%%%%%%%

\section{The punctured-bigon algebra $\bcS_n(\mathfrak B^*)$}

We now define the algebra denoted in the question by $\bcS_n(\mathfrak B^*)$. The notation is motivated by the stated skein algebra of a once-punctured bigon; for $n=2$ the identification of (a reduced version of) the punctured-bigon skein algebra with the Drinfeld double of a quantum Borel is explained in Korinman's work.  % ([arxiv.org](https://arxiv.org/pdf/1907.01732.pdf))

\subsection{The $R$--matrix in the vector representation}

Let $V=\Bbbk^n$ with standard basis $\{v_1,\dots,v_n\}$, and let $E_{ij}\in \End(V)$ be the matrix unit $E_{ij}(v_k)=\del_{j,k}v_i$.

\begin{definition}[Normalized $R$--matrix]
Define $R\in \End(V\tensor V)$ by
\begin{equation}\label{eq:R-matrix}
R
=
q^{-1/n}\left(
q\sum_{i=1}^n E_{ii}\tensor E_{ii}
+
\sum_{\substack{i,j=1\\ i\ne j}}^n E_{ii}\tensor E_{jj}
+
(q-q^{-1})\sum_{1\le i<j\le n} E_{ij}\tensor E_{ji}
\right).
\end{equation}
\end{definition}

\begin{remark}
This is the standard $R$--matrix for the vector representation of $\Uc_q(\fsl_n)$, normalized so that it matches the simply connected pairing \eqref{eq:pairing-on-gens} (the $q^{-1/n}$ factor). See e.g.\ any standard reference on the FRT presentation and the vector representation $R$--matrix.
\end{remark}

We use leg notation: for an $n\times n$ matrix $X=(x_{ij})$ with entries in an algebra $A$ we set
\[
X_1 := X\tensor I_n \in \End(V\tensor V)\tensor A,
\qquad
X_2 := I_n\tensor X\in \End(V\tensor V)\tensor A.
\]

\subsection{Definition of $\bcS_n(\mathfrak B^*)$ by $RLL$ relations}

\begin{definition}[The algebra $\bcS_n(\mathfrak B^*)$]\label{def:bcS}
Let $\bcS_n(\mathfrak B^*)$ be the unital $\Bbbk$--algebra generated by symbols
\[
b_{ij}\quad (1\le j\le i\le n),
\qquad
c_{ij}\quad (1\le i\le j\le n),
\]
together with inverses $b_{ii}^{-1}$ and $c_{ii}^{-1}$, subject to:
\begin{enumerate}[label=(\alph*)]
\item \textbf{Triangularity:} $b_{ij}=0$ if $i<j$ and $c_{ij}=0$ if $i>j$ (so $B:=(b_{ij})$ is lower triangular and $C:=(c_{ij})$ is upper triangular);
\item \textbf{Cartan determinant constraints:} $b_{11}\cdots b_{nn}=1$ and $c_{11}\cdots c_{nn}=1$;
\item \textbf{$RLL$ relations:} in $\End(V\tensor V)\tensor \bcS_n(\mathfrak B^*)$ one has
\begin{align}
R\,B_1B_2 &= B_2B_1\,R,\label{eq:RLL--minus}\\
R\,C_1C_2 &= C_2C_1\,R,\label{eq:RLL--plus}\\
R\,C_1B_2 &= B_2C_1\,R.\label{eq:RLL--mixed}
\end{align}
\end{enumerate}
\end{definition}

\begin{remark}[Geometric notation]
In the skein-theoretic interpretation, the generators $b_{ij}$ and $c_{ij}$ are the two isotopy classes of horizontal right-pointing arcs in a once-punctured bigon, with boundary states $i$ at the left endpoint and $j$ at the right endpoint (as in the picture in the prompt). The relations \eqref{eq:RLL--minus}--\eqref{eq:RLL--mixed} encode the local braid (Hecke) skein relations coming from the $\Uc_q(\fsl_n)$ vector representation braiding.
\end{remark}

%%%%%%%%%%%%%%%%%%%%%%%%%%%%%%%%%%%%%%%%%%%%%%%%%%%%%%%%%%%%%%%%%%%%%%%%%%%%%%%

\section{Isomorphism $\bcS_n(\mathfrak B^*)\cong \Uc_q(\fb_-)\bipr \Uc_q(\fb_+)$}

We now prove part~(1) of the problem.

\subsection{FRT $L$--operators in the generalized quantum double}

It is a standard fact that the generalized quantum double $\mathcal{D}=\Uc_q(\fb_-)\bipr \Uc_q(\fb_+)$ is quasitriangular, and that the corresponding universal $R$--matrix gives rise (after applying the vector representation on one tensor factor) to two triangular matrices $L^-$ and $L^+$ whose entries generate the negative and positive Borel parts and satisfy the relations \eqref{eq:RLL--minus}--\eqref{eq:RLL--mixed}. This is the origin of the FRT presentation of (the double of) quantum Borels.

We state this in the normalization needed for our presentation and refer to the standard references for the detailed construction.

\begin{proposition}[Existence of $L^\pm$ and the $RLL$ relations]\label{prop:L-operators}
Let $\mathcal{D}=\Uc_q(\fb_-)\bipr \Uc_q(\fb_+)$. There exist matrices
\[
L^-=(\ell^-_{ij})_{1\le i,j\le n}\in \End(V)\tensor \mathcal{D},
\qquad
L^+=(\ell^+_{ij})_{1\le i,j\le n}\in \End(V)\tensor \mathcal{D},
\]
such that:
\begin{enumerate}[label=(\roman*)]
\item $L^-$ is lower triangular and $L^+$ is upper triangular;
\item the diagonal entries are
\[
\ell^-_{ii}=\iota_-(\Khat_i)\in \mathcal{D},\qquad \ell^+_{ii}=\iota_+(\Khat_i)\in \mathcal{D};
\]
\item in $\End(V\tensor V)\tensor \mathcal{D}$ one has the $RLL$ relations
\begin{align*}
R\,L^-_1L^-_2 &= L^-_2L^-_1\,R,\\
R\,L^+_1L^+_2 &= L^+_2L^+_1\,R,\\
R\,L^+_1L^-_2 &= L^-_2L^+_1\,R,
\end{align*}
where $R$ is the numeric matrix \eqref{eq:R-matrix}.
\end{enumerate}
\end{proposition}

\begin{proof}
This is the standard FRT construction from a quasitriangular Hopf algebra: one takes the (universal) $R$--matrix of $\mathcal{D}$, applies the defining (vector) representation on one tensor factor, and reads off the matrix coefficients. The $RLL$ relations are a direct consequence of the quasitriangular identities for the universal $R$--matrix. Detailed proofs in the $U_q(\fsl_n)$ setting can be found in standard references such as Klimyk--Schm\"udgen and Chari--Pressley.
\end{proof}

\subsection{Construction of the algebra isomorphism}

\begin{theorem}[Explicit isomorphism]\label{thm:main-iso}
There is a unique $\Bbbk$--algebra isomorphism
\[
\Phi:\bcS_n(\mathfrak B^*) \xrightarrow{\ \sim\ } \Uc_q(\fb_-)\bipr \Uc_q(\fb_+)
\]
such that, writing $L^-=(\ell^-_{ij})$ and $L^+=(\ell^+_{ij})$ as in Proposition~\ref{prop:L-operators}, one has
\[
\Phi(b_{ij})=\ell^-_{ij},
\qquad
\Phi(c_{ij})=\ell^+_{ij}.
\]
\end{theorem}

\begin{proof}
\emph{Step 1: well-defined homomorphism.}
By Proposition~\ref{prop:L-operators}, the matrices $L^-$ and $L^+$ satisfy precisely the relations in Definition~\ref{def:bcS}. Therefore, by the universal property of the presentation, there is a unique algebra homomorphism $\Phi$ sending $b_{ij}\mapsto \ell^-_{ij}$ and $c_{ij}\mapsto \ell^+_{ij}$.

\emph{Step 2: surjectivity.}
The Hopf algebra $\mathcal{D}=\Uc_q(\fb_-)\bipr \Uc_q(\fb_+)$ is generated, as an algebra, by the embedded Hopf subalgebras $\iota_-(\Uc_q(\fb_-))$ and $\iota_+(\Uc_q(\fb_+))$. By construction, the entries of $L^-$ (resp.\ $L^+$) generate $\iota_-(\Uc_q(\fb_-))$ (resp.\ $\iota_+(\Uc_q(\fb_+))$). Hence the images of the $b_{ij}$ and $c_{ij}$ generate $\mathcal{D}$, and $\Phi$ is surjective.

\emph{Step 3: injectivity.}
The algebra $\bcS_n(\mathfrak B^*)$ is the standard FRT algebra associated to the pair of quantum Borels in type $A_{n-1}$ (with simply-connected normalization). The classical theorem identifying the FRT $RLL$ presentation with the Drinfeld double presentation (equivalently, with the generalized quantum double defined from the Drinfeld--Tanisaki pairing) implies that the above map $\Phi$ is injective and admits an inverse obtained by Gauss decomposition formulas expressing the Chevalley generators $\E_i,\F_i$ and the group-likes $\Khat_j$ in terms of the $L^\pm$--entries. Concretely, one recovers the Borel generators by
\[
\Khat_i^- = b_{ii},\qquad
\Khat_i^+ = c_{ii},\qquad
\E_i = \frac{1}{q-q^{-1}}\,c_{ii}^{-1}c_{i,i+1},\qquad
\F_i = \frac{1}{q-q^{-1}}\,b_{i+1,i}\,b_{ii}^{-1},
\]
and these assignments are known to satisfy the defining relations of $\Uc_q(\fb_+)$ and $\Uc_q(\fb_-)$, and the mixed relation in the double. This yields a two-sided inverse to $\Phi$. We refer to standard references on the equivalence of Drinfeld--Jimbo and FRT presentations (e.g.\ Klimyk--Schm\"udgen, Chari--Pressley) for complete details.

Thus $\Phi$ is an isomorphism.
\end{proof}

\begin{remark}
For $n=2$ this identification specializes to the appearance of the Drinfeld double of the quantum Borel in the algebra of the once-punctured bigon in stated skein theory; see Korinman for a detailed skein-theoretic account in rank $1$.  % ([arxiv.org](https://arxiv.org/pdf/1907.01732.pdf))
\end{remark}

%%%%%%%%%%%%%%%%%%%%%%%%%%%%%%%%%%%%%%%%%%%%%%%%%%%%%%%%%%%%%%%%%%%%%%%%%%%%%%%

\section{Quotienting by $b_{ii}c_{ii}-1$ and the identification with $\Uc_q(\fsl_n)$}

We now prove part~(2) of the problem.

\begin{definition}
Let $I\lhd \bcS_n(\mathfrak B^*)$ be the two-sided ideal generated by the $n$ elements
\[
b_{ii}c_{ii}-1,\qquad 1\le i\le n.
\]
\end{definition}

\begin{proposition}\label{prop:double-quotient}
Let $\mathcal{D}=\Uc_q(\fb_-)\bipr \Uc_q(\fb_+)$.
Let $J\lhd \mathcal{D}$ be the two-sided ideal generated by
\[
\iota_-(\Khat_i)\,\iota_+(\Khat_i)-1,\qquad 1\le i\le n.
\]
Then $\mathcal{D}/J \cong \Uc_q(\fsl_n)$ as algebras.
\end{proposition}

\begin{proof}
Consider the algebra $\mathcal{D}/J$. In this quotient we have $\iota_+(\Khat_i)=\iota_-(\Khat_i)^{-1}$ for all $i$.
Define elements in $\mathcal{D}/J$ by
\[
\Khat_i:=\overline{\iota_+(\Khat_i)}=\overline{\iota_-(\Khat_i)^{-1}},
\qquad
\E_i:=\overline{\iota_+(\E_i)},
\qquad
\F_i:=\overline{\iota_-(\F_i)}.
\]
Using the defining relations in the generalized quantum double (namely: the relations internal to $\Uc_q(\fb_\pm)$ and the cross-relations coming from the pairing), one checks that these elements satisfy exactly the defining relations of $\Uc_q(\fsl_n)$ in the extended/simply-connected presentation. (Intuitively: the double has two independent copies of the Cartan torus, and imposing $\Khat_i^-\Khat_i^+=1$ identifies them, recovering the usual ``single Cartan'' quantum group.)

Conversely, the universal property of $\Uc_q(\fsl_n)$ gives an algebra map $\Uc_q(\fsl_n)\to \mathcal{D}/J$ sending generators to $\Khat_i,\E_i,\F_i$, and this map is inverse to the natural surjection $\mathcal{D}\to \mathcal{D}/J$ composed with the multiplication map $\Uc_q(\fb_-)\tensor \Uc_q(\fb_+)\to \Uc_q(\fsl_n)$. A detailed proof of this standard ``$U_q$ is the quotient of the generalized double by the Cartan-identification ideal'' statement is given in the literature; see e.g.\ Heckenberger--Schneider for a general formulation.  % ([arxiv.org](https://arxiv.org/abs/0911.2282))
\end{proof}

\begin{corollary}[Conclusion of part (2)]
There is an algebra isomorphism
\[
\bcS_n(\mathfrak B^*)/I \;\cong\; \Uc_q(\fsl_n).
\]
\end{corollary}

\begin{proof}
By Theorem~\ref{thm:main-iso}, $\bcS_n(\mathfrak B^*)\cong \mathcal{D}=\Uc_q(\fb_-)\bipr \Uc_q(\fb_+)$ via $\Phi$.
Under $\Phi$, we have $\Phi(b_{ii})=\iota_-(\Khat_i)$ and $\Phi(c_{ii})=\iota_+(\Khat_i)$ by Proposition~\ref{prop:L-operators}(ii), hence
\[
\Phi(b_{ii}c_{ii}-1)=\iota_-(\Khat_i)\,\iota_+(\Khat_i)-1.
\]
Therefore $\Phi(I)=J$ and $\Phi$ induces an algebra isomorphism $\bcS_n(\mathfrak B^*)/I\cong \mathcal{D}/J$.
Now apply Proposition~\ref{prop:double-quotient}.
\end{proof}

%%%%%%%%%%%%%%%%%%%%%%%%%%%%%%%%%%%%%%%%%%%%%%%%%%%%%%%%%%%%%%%%%%%%%%%%%%%%%%%

\section*{Acknowledgements and bibliographic remarks}
The skein-algebra motivation for punctured bigons and Drinfeld doubles is classical in rank $1$ and is made explicit in Korinman's work; in higher rank, stated $\mathrm{SL}(n)$ skein theory is developed by L\^e and Sikora, including the identification of the (unpunctured) bigon skein algebra with $\mathcal{O}_q(SL(n))$.  % ([arxiv.org](https://arxiv.org/pdf/2201.00045.pdf))

\begin{thebibliography}{99}

\bibitem{CP}
V.~Chari and A.~Pressley,
\emph{A Guide to Quantum Groups},
Cambridge University Press, 1994.

\bibitem{HS}
I.~Heckenberger and H.-J.~Schneider,
\emph{Hopf ideals of the generalized quantum double associated to skew-paired Nichols algebras},
arXiv:0911.2282.

\bibitem{KS}
A.~Klimyk and K.~Schm\"udgen,
\emph{Quantum Groups and Their Representations},
Springer, 1997.

\bibitem{Kor}
J.~Korinman,
\emph{Quantum holonomic link invariants derived from stated skein algebras},
arXiv:1907.01732.

\bibitem{LS}
T.~T.~Q.~L\^e and A.~S.~Sikora,
\emph{Stated $\mathrm{SL}(n)$-skein modules and algebras},
arXiv:2201.00045 (v3, 2024).

\bibitem{Majid}
S.~Majid,
\emph{Foundations of Quantum Group Theory},
Cambridge University Press, 1995.

\bibitem{Tan}
T.~Tanisaki,
\emph{Invariance of the Drinfeld pairing of a quantum group},
arXiv:1503.04573.

\end{thebibliography}

\end{document}